%==============================================================================
\section{Skewness and Kurtosis --- Độ lệch và độ nhọn}
\label{sec:skewness-kurtosis}
%==============================================================================

\begin{dinhnghia}[title={Skewness (độ lệch)}]
Skewness đo \textbf{độ bất đối xứng} của phân phối.
\begin{itemize}
  \item Dương: đuôi dài bên phải.
  \item Âm: đuôi dài bên trái.
  \item Bằng 0: đối xứng (như phân phối chuẩn lý tưởng).
\end{itemize}
\end{dinhnghia}

\begin{dinhnghia}[title={Kurtosis (độ nhọn)}]
Kurtosis đo \textbf{độ nhọn} của phân phối so với chuẩn.
\begin{itemize}
  \item Kurtosis cao: đỉnh nhọn hơn, đuôi nặng hơn chuẩn.
  \item Kurtosis thấp: đỉnh dẹt hơn, đuôi nhẹ hơn.
  \item Bằng 0: gần chuẩn.
\end{itemize}
\end{dinhnghia}

\begin{ynghia}[title={Ảnh hưởng tới ML}]
Phân phối lệch mạnh có thể cần biến đổi dữ liệu hoặc phương pháp phi tham số.
Kurtosis cao có thể cần mô hình hoặc ước lượng robust hơn.
\end{ynghia}

\subsection{Ví dụ Python (SciPy)}

\begin{vidu}[title={skew() và kurtosis()}]
\begin{lstlisting}
import numpy as np
from scipy.stats import skew, kurtosis

data = np.random.normal(0, 1, 1000)

skewness = skew(data)
kurt_val = kurtosis(data)

print('Skewness:', skewness)
print('Kurtosis:', kurt_val)
\end{lstlisting}

\textbf{Output} (ví dụ, gần 0 với mẫu chuẩn):
\begin{lstlisting}
Skewness: -0.04119418903611285
Kurtosis: -0.1152250196054534
\end{lstlisting}
\end{vidu}
